\section{Positioning and physical interpretation}

The principal extension in this phase is a quantitative treatment of a
coarse detector whose natural entrance distribution need not coincide with the
reference preparation. A departure from ideal occupation ordering then has a
measurable consequence: any admissible classical explanation must provide enough
entrance mismatch, sufficiently rapid crossings of the observed boundary, or
both. The comparison is against a specified set of physical capabilities, rather
than against a fitted list of hidden networks.

Let $B$ be the detected region, let the reference preparation lie in $B$, and let
the comparison preparation lie outside it. Write $H(T)$ for the difference of
their expected cumulative occupations in $B$. In the reset-conditioned theorem,
$\eta$ bounds the total-variation mismatch between the first-entrance law and
the reference reset law, conditional on the entrance time. With the stated
uniform entry and exit intensity bounds $a,b$, define
\[
 A_c(T)=\int_0^T e^{-ct}\,dt,\qquad
 G_{a,b}(T)=\int_0^T e^{-bs-a(T-s)}\,ds.
\]
The sharp envelope is
\[
 H(T)\ge G_{a,b}(T)
 -\eta\{T-A_a(T)-A_b(T)+G_{a,b}(T)\}.
\]
Sharpness means attainability within the declared mathematical class; it does
not imply that an arbitrary apparatus realizes an extremizer. If the braces
are positive, a calibrated upper confidence bound on $H(T)$ gives a necessary
lower bound on $\eta$ for any classical explanation satisfying the chosen rate
caps. It gives neither an upper bound on the actual mismatch nor a unique hidden
mechanism.

At $\eta=1$, the envelope reduces to
$H(T)\ge A_a(T)+A_b(T)-T$. This endpoint requires only the relevant survival
envelopes and remains valid for history-dependent and time-dependent classical
processes satisfying the conditional hazard bounds. The intermediate
$0\le\eta<1$ theorem additionally uses its stated reset and continuation
conditions; the broader endpoint scope must not be transferred to it without
those conditions. Here $a$ and $b$ are caps on crossing intensities. They are not
the average total jump activity, heat production, or energy expenditure.

Classical speed and no-jump bounds are established. Shiraishi, Funo and Saito
relate distributional motion to dissipation and dynamical activity
\cite{rrShiraishi2018}; Hasegawa and Nishiyama derive concentration and response
bounds using no-jump probabilities \cite{rrHasegawa2024}. The present inverse
problem concerns a difference between two occupation experiments, with an
explicit reset-defect parameter and a saturating construction. Coarse-grained
memory and its effect on inference are also established concerns
\cite{rrBlom2024}. No claim of historical priority follows from these
comparisons. The contribution offered for review is this particular sharp
envelope, its operational interpretation, and the calibrated implementation.

Physical use requires evidence for the caps over the histories and preparations
covered by the test. An observed mean switching rate is insufficient. A hidden
entrance distribution cannot generally be certified from a binary detector
alone; when independent reset information is unavailable, the $\eta=1$ envelope
is the applicable conservative alternative. Preparation errors, detector drift,
timing discretization and finite-sample uncertainty must be included before
inverting the bound. A violation excludes the calibrated class. It does not by
itself establish quantum dynamics or identify the failed assumption. No
apparatus-specific validation of these requirements has yet been performed.

\subsection{Dimension rather than an assumed reset interface}

A complementary interpretation replaces reset and rate information with a
bound on hidden state count. Let $C_0$ denote the contrast between two arbitrary
preparations measured immediately, and $C_T$ their contrast with the same binary
detector at an independent uniform time in $[0,T]$. Every common homogeneous
classical Markov realization with at most three states obeys
\[
 |C_0-C_T|\leq1.
\]
The detector may have any fixed state-dependent response probability in $[0,1]$.
Measuring the immediate contrast removes the need to assume perfect extremal
preparations. The protected-sign rule follows when $C_0=1$. A violating qubit
example and a four-state classical construction distinguish quantum
attainability from the minimum classical realization size at this scalar
statistic. They do not characterize the full quantum record or exclude
classical explanations of unrestricted size.

Dynamical dimension inference already has a substantial literature. In
particular, delayed-data and spectral methods establish classical/quantum
dimension separations \cite{rrWolf2009}, while repeated-operation null tests
provide experimentally implemented alternatives \cite{rrBialecki2024}.
Temporal memory witnesses also predate this work
\cite{rrBudroni2019,rrVieira2022}. The proposed distinction is a finite-window
integrated inequality with an exact three-state threshold and a four-state
construction, under its stated preparation and evolution assumptions. Historical
priority and an experimental advantage over these methods remain unresolved.

Those assumptions supply real information. With two unrestricted preparations,
general prepare-and-measure data can always use two classical messages
\cite{rrGallego2010}; the stronger state count here depends on requiring a
single homogeneous continuous-time evolution and a common detector. Hidden
controller or environmental states that affect the evolution belong in the
classical realization being counted.
